% Quill starter using OUP's unmodified LPPL authoring class.
% Check Bioinformatics author guidelines before submission.
\documentclass[unnumsec,webpdf,contemporary,large,namedate]{oup-authoring-template}
\usepackage{originwrite-oup-xetex}
\newcommand{\societylogo}{}
\begin{document}
\journaltitle{Bioinformatics}
\DOI{Manuscript draft}
\copyrightyear{2026}
\pubyear{2026}
\access{Draft for review}
\appnotes{Original Paper}
\firstpage{1}
\title[Short title]{Your manuscript title}
\author[1]{Your name}
\address[1]{\orgdiv{Department}, \orgname{Institution}, \orgaddress{\country{Country}}}
\abstract{\textbf{Motivation:} Describe the problem.\\
\textbf{Results:} Summarise the approach and findings.\\
\textbf{Availability and implementation:} Give the software and data locations.}
\keywords{keyword one, keyword two}
\maketitle
\section{Introduction}
Explain the problem and its context.
\section{Methods}
Describe the approach and evaluation.
\section{Results}
Present the findings.
\section{Discussion}
Discuss limitations and implications.
\section{Data availability}
State where the supporting data and code can be accessed.
\section{Funding}
Describe funding and conflicts of interest.
% Link a personal library and sync references.bib before citing its papers.
% Uncomment these lines after inserting the first citation.
% \bibliographystyle{oup-abbrvnat}
% \bibliography{references}
\end{document}
